\acte{La Flèche d’or}
\usage{quot., trois fois ; dim., jour du Seigneur, et chaque fois qu’on entend un blasphème}
\rubrique{Acte de louange pour la réparation des blasphèmes du saint Nom de Dieu, dictée par Notre-Seigneur à la sœur Marie de Saint-Pierre.}
\ouverture{Q}{u’à}{jamais} soit loué, béni, aimé, adoré, glorifié le très saint, très sacré, très adorable, très inconnu, très inexprimable Nom de Dieu, au ciel, sur la terre et dans les enfers, par toutes les créatures sorties des mains de Dieu et par le Sacré-Cœur de Notre-Seigneur Jésus-Christ au très saint Sacrement de l’autel !
\conclusion{Ainsi soit-il.}
\rubrique{On dira trois fois cet acte de louange en l’honneur des trois personnes de la sainte Trinité.}
\begin{notice}
Dans les écrits de la sœur on lit :

« Notre-Seigneur m’a ouvert son Cœur, et j’ai entendu ces paroles : Mon Nom est partout blasphémé ; même les enfants blasphèment. Alors il m’a fait entendre combien cet affreux péché blessait son divin Cœur douloureusement et plus que tous les autres ; par le blasphème, le pécheur le maudit en face, l’attaque ouvertement, anéantit la rédemption et prononce lui-même sa condamnation et son jugement. Le blasphème est comme une flèche empoisonnée qui blesse continuellement son Cœur ; et il me dit qu’il voulait me donner une flèche d’or pour le blesser délicieusement et cicatriser les blessures de malice que lui font les pécheurs.

« Voici la formule de louange qu’il me dicta, malgré ma grande indignité, pour la réparation des blasphèmes de son saint Nom, et qu’il me donna comme une flèche d’or, m’assurant qu’à chaque fois que je la dirai je blesserai son Cœur d’une blessure d’amour.

« Il ajouta : Faites attention à cette faveur ; car je vous en demanderai compte. À ce moment il me semblait voir sortir du Sacré-Cœur de Jésus blessé par cette flèche d’or des torrents de grâces pour la conversion des pécheurs. »
\end{notice}
\source{Dévotion à la Sainte Face de Notre-Seigneur Jésus-Christ, 1887}
\culdelampe